\documentclass{article}
\usepackage[utf8]{inputenc}
\usepackage[T1]{fontenc}
\usepackage{geometry}
\usepackage{graphicx}
\usepackage{float}
\usepackage{hyperref}
\usepackage[english,indonesian]{babel}

\title{Optimasi Pemrosesan Template LaTeX Menggunakan Pendekatan Berbasis AI}
\author{Academic Support Agent}
\date{\today}

\begin{document}
\maketitle

\begin{abstract}

\end{abstract}

\chapter{PENDAHULUAN}


\chapter{TINJAUAN PUSTAKA}


% Original content (commented out):
% \documentclass[letterpaper]{article}
\usepackage{amsmath,amssymb}
\usepackage{lmodern}
\usepackage{iftex}

\usepackage[T1]{fontenc}
\usepackage[utf8]{vietnam}
\usepackage{textcomp} % provide euro and other symbols
\usepackage{lipsum}
\usepackage{setspace}
\usepackage{hyperref}
\usepackage{xcolor}
\usepackage{longtable,booktabs,array}
\usepackage{multirow}
\usepackage{calc} % for calculating minipage widths
% Correct order of tables after \paragraph or \subparagraph
\usepackage{etoolbox}

\usepackage{adjustbox}
\usepackage{graphicx}
\makeatletter
\def\maxwidth{\ifdim\Gin@nat@width>\linewidth\linewidth\else\Gin@nat@width\fi}
\def\maxheight{\ifdim\Gin@nat@height>\textheight\textheight\else\Gin@nat@height\fi}
\makeatother
% Scale images if necessary, so that they will not overflow the page
% margins by default, and it is still possible to overwrite the defaults
% using explicit options in \includegraphics[width, height, ...]{}
\setkeys{Gin}{width=\maxwidth,height=\maxheight,keepaspectratio}

\usepackage{tabularx}
% \setlength{\parskip}{1cm plus 4mm minus 3mm}
% \setlength{\parindent}{1cm}
\usepackage[top=15mm, bottom=15mm, left=30mm, right=30mm]{geometry} % page layout

% add biblatex
\usepackage[backend=bibtex,style=numeric ,natbib=true, sorting=nty, maxbibnames=99]{biblatex} 
\addbibresource{ref}

% re-define section/ subsection
\usepackage{titlesec}
\providecommand{\tightlist}{%
  \setlength{\itemsep}{0pt}\setlength{\parskip}{0pt}}
\setcounter{secnumdepth}{-\maxdimen} % remove section numbering
\DeclareUnicodeCharacter{0301}{*************************************}

\titleformat{\section}
  {\large\bfseries\color{blue}}
  {\thesection}{1em}{}
\titleformat{\subsection}
  {\large\bfseries\color{blue}}
  {\thesection}{1em}{}
\titleformat{\subsubsection}
  {\normalsize\color{blue}}
  {\thesection}{1em}{}
\author{Khoa Nguyen}
\date{Oct 2023}

\begin{document}

\begin{table*}[!ht]
    \centering
    \begin{tabular}{m{6em} c  m{7em} m{5em} }
   

         \multirow{3}{3em}{
         \begin{minipage}{0.15\textwidth}
      \includegraphics{asset/Figure/image1.png}
    \end{minipage}} & 
     & & \\  \cline{3-4}
     & \small ĐẠI HỌC QUỐC GIA TP.HCM &  \multicolumn{1}{|c|}{\multirow{2}{8em}{\small Ngày nhận hồ sơ}} &  \multicolumn{1}{|c|}{}\\
    & \small \textbf{TRƯỜNG ĐẠI HỌC CÔNG NGHỆ THÔNG TIN} & \multicolumn{1}{|c|}{} &  \multicolumn{1}{|c|}{}\\\cline{3-4}

        & & \multicolumn{2}{|c|}{\textit{\small (Do CQ quản lý ghi)}}\\ \cline{3-4}
    
    \end{tabular}
    % \caption{Caption}
    % 
\end{table*}

\begin{quote}
\centering
\textbf{\Large THUYẾT MINH}

\Large ĐỀ TÀI KHOA HỌC VÀ CÔNG NGHỆ CẤP SINH VIÊN 2023
\end{quote}


\section{A. THÔNG TIN CHUNG}

\subsection{A1. Tên đề tài}


\begin{itemize}
\item[-]  Tên tiếng Việt (IN HOA): 
\item[-]  Tên tiếng Anh (IN HOA): 
\end{itemize}


\subsection{A2. Thời gian thực hiện}

\textbf{\ldots} tháng (kể từ khi được duyệt).

\subsection{A3. Tổng kinh phí}

\emph{(Lưu ý tính nhất quán giữa mục này và mục B8. Tổng hợp kinh phí đề
nghị cấp)}

\noindent Tổng kinh phí: \textbf{\ldots} 
triệu đồng, gồm:

\begin{itemize}
\item
  Kinh phí từ Trường Đại học Công nghệ Thông tin: \textbf{\ldots} triệu
  đồng
\end{itemize}


\subsection{A4. Chủ nhiệm}
\begin{spacing}{1.5}

\indent \indent Họ và tên: \par

Ngày, tháng, năm sinh: \hspace{2cm} Giới tính (Nam/Nữ): \par

Số CMND: \hspace{3cm}; Ngày cấp: \hspace{4cm}; Nơi cấp: 

Mã số sinh viên:  

Số điện thoại liên lạc: 

Đơn vị (Khoa): Kỹ thuật máy tính 

Số tài khoản: \hspace{4cm} Ngân hàng: 
\end{spacing}

\subsection{A5. Thành viên đề tài}


\begin{longtable}[c]{@{}
  >{\raggedright\arraybackslash}p{(\columnwidth - 6\tabcolsep) * \real{0.0553}}
  >{\raggedright\arraybackslash}p{(\columnwidth - 6\tabcolsep) * \real{0.3235}}
  >{\raggedright\arraybackslash}p{(\columnwidth - 6\tabcolsep) * \real{0.3637}}
  >{\raggedright\arraybackslash}p{(\columnwidth - 6\tabcolsep) * \real{0.2575}}@{}}
\toprule()
\begin{minipage}[b]{\linewidth}\raggedright
\textbf{TT}
\end{minipage} & \begin{minipage}[b]{\linewidth}\raggedright
\textbf{Họ tên}
\end{minipage} & \begin{minipage}[b]{\linewidth}\raggedright
\textbf{MSSV}
\end{minipage} & \begin{minipage}[b]{\linewidth}\raggedright
\textbf{Khoa}
\end{minipage} \\
\midrule()
\endhead
1 & Nguyễn Văn A & 18521234 & Khoa học Máy tính \\
\bottomrule()
\end{longtable}




\section{B. MÔ TẢ NGHIÊN CỨU}
\indent \indent Nội dung mô tả nghiên cứu được trình bày ở đây. Có thể tham chiếu đến các kết quả nghiên cứu khác thông qua \textbf{ \textbackslash cite\{tham chiếu paper trong file .bib\}} như ví dụ này \cite{vaswani2017attention}.


Ngoài ra, bạn có thể tạo ảnh minh họa với từ khóa \textbf{ \textbackslash begin\{figure\}...\textbackslash end\{figure\}} và tham chiếu đến ảnh với từ khóa \textbf{ \textbackslash ref\{<label của ảnh>\}} giống ví dụ sau \ref{fig:trans-arch}.

\begin{figure}[!hbt]
    \centering
    \includegraphics[width=0.7\textwidth]{Figures/transformer_arch.png}
    \caption{Kiến trúc của mô hình Transformer~\cite{vaswani2017attention}}
    \label{fig:trans-arch}
\end{figure}



\subsection{B1. Giới thiệu về đề
tài}

\subsection{B2. Mục tiêu, nội dung, kế hoạch nghiên cứu}
\subsection{B2.1 Mục tiêu}
\subsection{B2.2 Nội dung và phương pháp nghiên cứu}

\subsubsection{Nội dung 1: Tìm hiểu tổng quan đề tài. }

\paragraph{Phương pháp thực hiện:}

\paragraph{Kết quả dự kiến:}

\subsubsection{Nội dung 2: Đề xuất độ đo đặc trưng hình ảnh cho bài toán nhận diện khuôn mặt.}

\paragraph{Phương pháp thực hiện:}

\paragraph{Kết quả dự kiến:}

\subsubsection{Nội dung 3: Hiện thực bài toán nhận diện khuôn mặt trên thiết
bị nhúng. }

\paragraph{Phương pháp nghiên cứu:}

\paragraph{Kết quả dự kiến:}

\subsubsection{Nội dung 4: Đóng gói sản phẩm và viết báo cáo.}

\paragraph{Phương pháp nghiên cứu:}

\paragraph{Kết quả dự kiến:}

\subsection{B2.3 Kế hoạch nghiên cứu}



\subsection{B3. Kết quả dự kiến}

\printbibliography[heading=bibintoc, title={B4. Tài liệu tham
khảo}]



\centering
\begin{longtable}[c]{@{}
  >{\centering\arraybackslash}p{(\columnwidth - 2\tabcolsep) * \real{0.6}}
  >{\centering\arraybackslash}p{(\columnwidth - 2\tabcolsep) * \real{0.4}}@{}}
% \toprule()
    \begin{minipage}[b]{\linewidth}\centering
        \emph{Ngày \ldots tháng \ldots năm 2023}\\
        \textbf{Giảng viên hướng dẫn}\\
        (Ký và ghi rõ họ tên)
    \end{minipage} & 
    \begin{minipage}[b]{\linewidth}\centering
        \emph{Ngày \ldots tháng \ldots năm 2023}\\
        \textbf{Chủ nhiệm đề tài}\\
        (Ký và ghi rõ họ tên)
    \end{minipage} \\
\end{longtable}

\end{document}




\end{document}